\chapter{Experiments and Results}\label{chap:resultats}

\section{Introduction}
\TODO{Chapter introduction}

\section{Experimental Environment}
% Table: Python, libraries and versions (requirements.txt), random seed, hardware.
\TODO{Software and hardware versions table}

\section{Cross-Validation Results}
% Comparative table of the six models: macro-F1, weighted F1, balanced accuracy, standard deviation across folds.
% Do not mix these values with test-set results in the same table.
\TODO{Results table and interpretation}

\section{Test Set Results}
% Selected model: accuracy, macro-F1, confusion matrix (figure), per-class report, one-vs-rest ROC-AUC.
% Explain the difference between cross-validation and test results.
\TODO{Results table, confusion matrices, and ROC curves}

\section{Rare Class and Data Leakage Features}
% Reminder about the reset-both class: report the actual counts, not just percentages.
% Ablation study with and without NAT-related features, if available; otherwise, honestly state this as a limitation.
\TODO{Rare-class analysis and ablation study}

\section{Discussion}
% What the results demonstrate, what they do not demonstrate, external validity (only one firewall).
\TODO{Discussion in prose}

\section{Conclusion}
\TODO{Chapter conclusion and transition to Chapter~\ref{chap:application}}